\subsection{Image and text (CC3M, CC12M)}
\label{sec:exp:imagetext}

\textbf{Setup.} We train CoBit-CC3M (462M parameters, $24\times1024$) on Conceptual Captions 3M
\citep{image:cc3m} ($\approx$2.9M pairs, $10^6$ steps) and CoBit-CC12M (633M, $26\times1152$) on
Conceptual 12M \citep{image:cc12m} ($\approx$11.7M pairs, $1.744\times10^6$ steps, 87\% of its
planned budget). Images are tokenized by a frozen Open-MAGVIT2 lookup-free quantiser
\citep{image:magvit2,image:lfq} into 18-bit codes, $16{\times}16=256$ tokens ($f16$) for CC3M
and $32{\times}32=1{,}024$ ($f8$) for CC12M; text by the \texttt{o200k} byte-pair vocabulary
(caption budgets of 48 and 96 tokens). A 19th namespace bit separates image from text and control codes, so every
position is one 19-bit token (\S\ref{sec:method:multimodal}); training regimes are $0.35$
text-conditioned, $0.35$ image-conditioned, $0.30$ joint. Neither model sees MS-COCO
\citep{image:mscoco}: all numbers are zero-shot on COCO \texttt{val2014} under one fixed protocol
(FID-30K, CMMD \citep{image:cmmd}, CLIP ViT-B/32 \citep{image:clip}; captioning on the Karpathy
test split; Appendix~\ref{app:image:protocol}). Deliberately absent are text-only pre-training, tokenizer surgery, auxiliary alignment losses and
staged curricula. All are known to help and we expect them to help here; we omit them so that
the result measures what the unmodified framework transfers, and leave them to future work.

\begin{table}[t]
\centering
\small
\setlength{\tabcolsep}{4.5pt}
\begin{tabular}{lrrrrrr}
\toprule
\textbf{Direction} & $n$ & Steps (NFE) & \textbf{FID}$\downarrow$ & \textbf{CMMD}$\downarrow$ & \textbf{CLIP}$\uparrow$ & \textbf{CIDEr}$\uparrow$ \\
\midrule
\multicolumn{7}{l}{\emph{CoBit-CC3M --- one set of weights (462M, $f16$, 256 image tokens)}}\\
text $\rightarrow$ image & 30{,}000 & 10 (20) & \textbf{18.48} & \textbf{18.41} & 27.73 & --- \\
image $\rightarrow$ text (CLIP) & 5{,}000 & 256 (512) & --- & --- & 24.83 & --- \\
image $\rightarrow$ text (captioning) & 5{,}000 & 256 (512) & --- & --- & --- & \textbf{0.136} \\
joint, 10 steps & 30{,}000 & 10 (10) & 50.19 & 41.53 & 23.10$^{\dagger}$ & --- \\
joint, 256 steps & 30{,}000 & 256 (256) & 61.95 & 46.80 & 25.67$^{\dagger}$ & --- \\
\midrule
\multicolumn{7}{l}{\emph{CoBit-CC12M --- one set of weights (633M, $f8$, 1{,}024 image tokens)}}\\
text $\rightarrow$ image & 30{,}000 & 16 (32) & \textbf{16.59} & \textbf{17.76} & 26.51 & --- \\
image $\rightarrow$ text (CLIP) & 5{,}000 & 256 (512) & --- & --- & 22.98 & --- \\
image $\rightarrow$ text (captioning) & 5{,}000 & 256 (512) & --- & --- & --- & 0.057 \\
joint, 20 steps & 5{,}000$^{\ddagger}$ & 20 (20) & 56.18 & 37.60 & 23.19$^{\dagger}$ & --- \\
joint, 256 steps & 5{,}000$^{\ddagger}$ & 256 (256) & 64.47 & 40.33 & 24.01$^{\dagger}$ & --- \\
\midrule
\emph{rFID floor of the tokenizer, $f16$ / $f8$} & 30{,}000 & & 1.59 / 0.73 & --- & --- & --- \\
\bottomrule
\end{tabular}
\caption{Zero-shot MS-COCO, all four generative directions from a single set of weights per
model. Each conditional row is quoted at the step count that is best for that direction
(\emph{Sampler} below); joint (unconditional) generation has no single optimum, so both ends of
its trade-off are shown. CLIP is ViT-B/32; per-direction guidance and churn are listed in
Appendix~\ref{app:image:protocol}, and ViT-L/14 scores, SWD and precision/recall in Appendix
Table~\ref{tab:app:image:headline-full}. $^{\dagger}$Pair
coherence $\cos(\text{generated image},\text{generated caption})$, not comparable to the
conditional CLIP scores; FID of the joint rows is of the generated image half. $^{\ddagger}$FID is
strongly biased in $n$ ($+5.3$ at $n{=}5000$ versus $30{,}000$), so these rows are not
FID-comparable to CC3M's (CC3M at $n{=}5000$: 56.06 / 67.47 at 10 / 256 steps).}
\label{tab:image:headline}
\end{table}

\textbf{Results.} Table~\ref{tab:image:headline} reports all four directions for both models,
each from a single set of weights. CC3M reaches zero-shot FID-30K \textbf{18.48} at NFE 20; its
tokenizer's own reconstruction floor on the same 30{,}000 images is rFID 1.59, so the codec
accounts for under 2 of those points and the gap to larger systems lies in the model, not in the
token space. Joint co-generation is much harder than either conditional direction (FID 50.19 at
10 steps), as with no prompt the model samples its own CC3M-like training distribution, yet the
co-generated image and caption agree with each other (pair coherence 25.67 at 256 steps) about
as well as a conditional caption agrees with its image (24.83).

\begin{table}[t]
\centering
\small
\setlength{\tabcolsep}{4pt}
\begin{tabular}{lrrrr}
\toprule
\textbf{Model} & \textbf{Params} & \textbf{Pre-train pairs} & \textbf{FID-30K}$\downarrow$ & \textbf{GenEval}$\uparrow$ \\
\midrule
\multicolumn{5}{l}{\emph{CoBit (ours; from scratch, no pretrained components)}}\\
CoBit-CC3M & 462M & 2.9M & \textbf{18.48} & 0.249$^{*}$ \\
CoBit-CC12M & 633M & 12M & \textbf{16.59} & --- \\
\midrule
\multicolumn{5}{l}{\emph{Same pre-training corpus (CC3M), zero-shot COCO}}\\
LAFITE \citep{image:lafite} & 75M ($+$CLIP) & 2.9M & 26.94 & --- \\
\midrule
\multicolumn{5}{l}{\emph{Unified image$+$text models}}\\
CoDI \citep{image:codi} & --- & 400M & 22.26 & 0.31 \\
SEED-X \citep{image:seedx} & 17B & --- & 14.99 & 0.49 \\
LWM \citep{image:lwm} & 7B & --- & 12.68 & 0.47 \\
UniDiffuser \citep{image:unidiffuser} & $\sim$1B & LAION-5B (subsets) & 9.71 & --- \\
Show-o \citep{image:showo} & 1.3B & 35M & 9.24 & 0.68 \\
\midrule
\multicolumn{5}{l}{\emph{Text-to-image specialists}}\\
DALL$\cdot$E \citep{image:dalle} & 12B & 250M & 27.50 & --- \\
CogView \citep{image:cogview} & 4B & 30M & 27.10 & --- \\
minDALL$\cdot$E & --- & --- & --- & 0.23 \\
LlamaGen & 0.8B & --- & --- & 0.32 \\
LDM \citep{image:ldm} & 1.4B & 400M & 12.64 & 0.37 \\
DALL$\cdot$E~2 \citep{image:dalle2} & 6.5B & 650M & 10.39 & 0.52 \\
SD v1.5 \citep{image:ldm} & 0.9B & 2000M & 9.62 & 0.43 \\
PixArt-$\alpha$ \citep{image:pixart} & 0.6B & 25M & 7.32 & 0.48 \\
Imagen \citep{image:imagen} & 3B & 860M & 7.27 & --- \\
\bottomrule
\end{tabular}
\caption{Zero-shot MS-COCO FID-30K and GenEval overall score against published systems. Prior
figures are as printed in the cited papers and in the compilation tables of \citet{image:showo}
(Tables 2--3), \citet{image:unidiffuser} (Table 1), \citet{image:lafite} (Table 2) and
\citet{image:geneval} (Table 2); parameter and pair counts as printed by \citet{image:showo}
where available. Rows are not strictly comparable (see text): corpora, backbones, pretrained
components and FID pipelines all differ, and the same nominal metric varies by about one FID
point between sources (Stable Diffusion is 8.59 in one compilation and 9.62 in another).
Improved VQ-Diffusion's 8.44 \citep{image:vqdiffusion-improved} is omitted because that model
is trained on COCO. $^{*}$Measured at 256 sampling steps, $25\times$ past the image optimum; a
lower bound.}
\label{tab:image:litcomparison}
\end{table}

\textbf{Comparison to prior work.} Table~\ref{tab:image:litcomparison} places the result among
published zero-shot MS-COCO numbers. Exact comparability is not attainable in this setting and we
do not claim it; the table is read for where the model sits, not for rank. Against the one
system with the same pre-training corpus and protocol, LAFITE, FID improves by 8.46 without its
frozen CLIP encoder. Among unified image$+$text models we sit between CoDI (400M pairs) and
SEED-X (17B parameters), and ahead of the DALL$\cdot$E and CogView specialists at 12B and 4B
parameters. Every system with a lower FID than ours relies on one or more of a text-pretrained
backbone or pretrained text encoder, a pretrained CLIP or VAE, and at least an order of magnitude
more paired data---precisely the ingredients withheld here.\note{Verify this sentence row by row
against the cited papers before submission.} On GenEval, 0.249 sits above minDALL$\cdot$E and
below CoDI and LlamaGen, far below the LLM-backed unified models. For captioning no zero-shot COCO baseline has
been published; COCO-trained specialists reach CIDEr 1.02--1.40 against our 0.136
(\emph{Limitations} below). Taken together, a 462M-parameter model trained from scratch on 2.9M
alt-text pairs, with no modality-specific machinery, generates images in the range of unified
systems one to two orders of magnitude larger and captions and co-generates from the same
weights. That is the claim: capable transfer, not a leaderboard entry.

\textbf{Sampler: the two modalities want step counts $\mathbf{25\times}$ apart.} Sweeping the
step count from 4 to 512 on the same weights (Appendix~\ref{app:image:sampler}),
text$\rightarrow$image FID and CMMD minimise at 10 steps for CC3M (16--20 for CC12M) and degrade
monotonically beyond; at $n{=}30{,}000$ the conventional 256 steps costs 4.06 FID for
$25.6\times$ the compute. Captioning inverts this: CIDEr is 0.013 at 10 steps, peaks at 0.136 at
256 and falls to 0.103 at 512 (image$\rightarrow$text CLIP behaves alike). The joint direction shows both effects inside single samples: every
step added improves the text half and degrades the image half (10, 32, 64 steps at $n{=}5000$:
pair-CLIP 23.16, 24.71, 25.17; FID 56.06, 59.14, 62.66), hence the two joint rows of
Table~\ref{tab:image:headline}. This is not a modification of the framework: the step count, like
the per-direction churn and guidance (Appendix~\ref{app:image:sampler}), is an inference-time
hyperparameter of any diffusion sampler, set per direction on the same weights without
retraining---and such sampler settings are the only thing that had to be adapted per modality.

\textbf{Scaling the corpus.} Moving to CC12M ($4\times$ the pairs, $1.37\times$ the parameters,
$\approx8.7\times$ the compute, a better codec) improves images and degrades text:
text$\rightarrow$image FID falls $18.48\rightarrow16.59$, a margin that survives subtracting each
tokenizer's own rFID floor ($16.89$ versus $15.87$; CMMD, $18.41\rightarrow17.76$, has no such
correction and retains an unquantified codec advantage), while at a matched 256 steps captioning
CIDEr falls $0.136\rightarrow0.057$ and image$\rightarrow$text CLIP $24.83\rightarrow22.98$. This
split is what noisier CC12M alt-text predicts and uniform undertraining does not, but that is an
inference from a pattern across metrics, not a measurement; the decisive checkpoint-progression
test has not been run (Appendix~\ref{app:image:scaling}).

\textbf{Limitations.} Compositional prompt following is near the floor: CC3M scores 0.249 on
GenEval \citep{image:geneval} (single object 0.80, colours 0.46, position 0.04, attribute
binding 0.03; Appendix~\ref{app:image:geneval}). This was measured at 256 steps, $25\times$ past
the image optimum, so it is a lower bound; CC12M has not been evaluated on GenEval. Captioning is weak in absolute terms (CIDEr 0.136, roughly $7.5\times$
below the COCO-trained L-Verse \citep{image:lverse} at 1.022; no zero-shot COCO captioning
baseline has been published), and the outputs are fluent CC3M-style alt-text scored against
COCO's terse references (RefCLIPScore \citep{image:clipscore} 0.661;
Appendix~\ref{app:image:captioning}). Both limitations are
consistent with CC3M alt-text rarely specifying relations, attribute bindings or COCO-style
descriptions, and both are exactly what the omitted ingredients---text pre-training, curated or
re-captioned data, alignment losses---are expected to buy. That is the point of the experiment:
it shows what unmodified transfer does and does not provide.
